%%%%%%%%%%%%%%%%%%%%%%%%%%%%%%%%%%%%%%%%%%%%%%%%%%%%%%%%%%%%%%%%%%%%%%
% Lớp 11 | Hàm số mũ và hàm số lôgarít | Phép tính lôgarít
% Dạng 1D6?2-3: Rút gọn, chứng minh biểu thức lôgarít
% Số câu: 82  (sinh tự động bởi tach_cau.py; từ nay thầy sửa trực tiếp file này)
%%%%%%%%%%%%%%%%%%%%%%%%%%%%%%%%%%%%%%%%%%%%%%%%%%%%%%%%%%%%%%%%%%%%%%

%%%%%%%%%%%%%%%% TRẮC NGHIỆM 4 PHƯƠNG ÁN %%%%%%%%%%%%%%%%

% nguồn: Ok_TracNghiemLop11_Lan1
\begin{ex}%[1D6N2-3]%[ma:8ddf8627]
	Với mọi số thực dương $a$, $b$, $x$, $y$  và $a$, $b\ne 1$, mệnh đề nào sau đây {\bf sai}?
	\choice 
		{\True  $\log_a\dfrac{1}{x}=\dfrac{1}{\log_ax}$}
		{$\log_a\left(xy\right)=\log_ax+\log_ay$} 
		 {$\log_ba \cdot \log_ax=\log_bx$}
		 {$\log_a\dfrac{x}{y}=\log_ax-\log_ay$}
	 \loigiai{
		Với mọi số thực dương $a$, $b$, $x$,  $y$ và $a$, $b\ne 1$, ta có $$\log_a\dfrac{1}{x}=\log_a{x^{-1}}\ne\dfrac{1}{\log_ax}.$$
		Vậy mệnh đề $\log_a\dfrac{1}{x}=\dfrac{1}{\log_ax}$ sai. Theo các tính chất logarit  thì các mệnh đề còn lại đều đúng.
	} 
\end{ex}

% nguồn: Ok_TracNghiemLop11_Lan1
\begin{ex}%[1D6N2-3]%[ma:04e7617c]
	Cho $a$ là một số thực dương và biểu thức $T= \log_4(4a)$. Khẳng định nào sau đây  đúng?
	\choice 
		{$T = 4\log_4 a$} 
		{$T =  \log_4 a$}
		{$T  = 1-\log_4 a$}
		{\True $T  = 1+\log_4 a$}
	\loigiai{ 
		Ta có  $T= \log_4(4a) = \log_4 4 + \log_4 a = 1+ \log_4 a$.
	} 
\end{ex}

% nguồn: Ok_TracNghiemLop11_Lan1
\begin{ex}%[1D6N2-3]%[ma:b54f7bfc]
	Cho $ a,b > 0$. Rút gọn biểu thức $\log_a{b^2}+\log_{a^2}{b^4}$  ta được
	 \choice
		{$ 2\log_ab$} 
		 {$ 0$}
		 {$\log_ab$}
		 {\True $ 4\log_ab$}
	\loigiai{ 
	Ta  có $\log_a{b^2}+\log_{a^2}{b^4}=2\log_ab+\dfrac{1}{2}\cdot 4\cdot \log_ab=4\log_ab$.}
\end{ex}

% nguồn: Ok_TracNghiemLop11_Lan1
\begin{ex}%[1D6N2-3]%[ma:f29d2b40]
	Cho $a, b, c > 0; a \neq 1$ và số $\alpha \in \mathbb{R}$, Khẳng định nào sau đây  \textbf{sai}?
	\choice 
		 {$\log_a a=1$}
		{\True $\log  (a+b)=\log a+\log b$}
		{$\log_a a^c=c$} 
		{$\log_a b^{\alpha}=\alpha \log_a b$} 
	\loigiai{ 
		Theo tính chất của logarit ta có khẳng định  sai là \lq\lq $\log (a+b)=\log a+\log b$ \rq\rq.
	} 
\end{ex}

% nguồn: Ok_TracNghiemLop11_Lan1
\begin{ex}%[1D6N2-3]%[ma:df7ac885]
	Với $a$, $b$ là các số thực dương và $m$, $ n$  là các số nguyên, mệnh đề nào sau đây \textbf{sai}?
	 \choice
		{$\log a+\log b=\log  a b$}
		{$\log a-\log  b=\log \dfrac{a}{b}$}
		{$\log_a  1=0$}
		{\True $\log a+\log b=\log a \cdot  \log b$}
	\loigiai{ 
		Mệnh đề sai là \lq\lq$\log a+\log b=\log  a \cdot \log b$\rq\rq.
	} 
\end{ex}

% nguồn: Ok_TracNghiemLop11_Lan1
\begin{ex}%[1D6N2-3]%[ma:64a984d8]
	Với hai số thực dương $a$, $b$ thỏa mãn $\log_2 a-3\log_2 b=2$,  khẳng định nào dưới đây đúng?
	 \choice
		 {$a=3b+2$}
		{$a=3b+4$} 
		{$a=\dfrac{4}{b^3}$} 
		 {\True $a=4b^3$}
	\loigiai{ 
	 Với hai số thực dương $a$, $b$ ta có $\log_2 a-3\log_2 b=2 \Leftrightarrow \log_2 \left(\dfrac{a}{b^3}\right)=2 \Leftrightarrow \dfrac{a}{b^3}=4 \Leftrightarrow a=4b^3$.
	} 
\end{ex}

% nguồn: Ok_Toan11_TracNghiem_Lan2
\begin{ex}%[1D6H2-3]%[ma:b20c229e]
	Cho các số thực dương $a$, $b$ thỏa mãn  $2\log_32\cdot \log_2a-3\log_{\sqrt 3}b=4$, mệnh đề nào dưới đây đúng?
	\choice 
		 {$a=\dfrac{9}{b^3}$}
		 {$a=4b$}
		 {$ab=4$}
		 {\True $a=9b^3$}
	 \loigiai{
		Ta  có
		 \allowdisplaybreaks
		\begin{eqnarray*}
			&&2\log_32\cdot  \log_2a-3\log_{\sqrt 3}b=4\\
			&\Leftrightarrow&2\log_3a-3\log_{3^{\tfrac{1}{2}}}b=4\\ 
			&\Leftrightarrow&\log_3a-3\log_3b=2\\ 
			&\Leftrightarrow&\log_3a-\log_3b^3=2\\ 
			 &\Leftrightarrow&\log_3\dfrac{a}{b^3}=2\\
			&\Leftrightarrow&\dfrac{a}{b^3}=3^2\\ 
			 &\Leftrightarrow&a=9b^3.
		\end{eqnarray*}
	 }
\end{ex}

% nguồn: Ok_Toan11_TracNghiem_Lan2
\begin{ex}%[1D6H2-3]%[ma:167dc313]
	Với $a$, $b$ là các số thực dương tùy  ý thỏa mãn $\log_22a^4+\log_2b=5$, khẳng định nào sau đây đúng?
	\choice 
		 {$a^4b=32$}
		 {$a^4+b=32$}
		{$a^4+b=16$} 
		{\True $a^4b=16$} 
	\loigiai{ 
		 Ta có $\log_22a^4+\log_2b=5\Leftrightarrow (2a^4b)=2^5\Leftrightarrow a^4b=2^4=16$.
	} 
\end{ex}

% nguồn: Ok_Toan11_TracNghiem_Lan2
\begin{ex}%[1D6H2-3]%[ma:a48ea822]
	Cho các số thực dương $a$, $b$ thỏa mãn $\log a+2\log  b=1$. Mệnh đề nào sau đây đúng?
	\choice 
		{$a+b^2=1$} 
		 {$a+2b=10$}
		 {\True $ab^2=10$}
		{$a+b^2=10$} 
	\loigiai{ 
		Ta có  \[\log a+2\log b=1\Leftrightarrow\log a+\log b^2=1\Leftrightarrow \log ab^2=1\Leftrightarrow ab^2=10.\]
	 }
\end{ex}

% nguồn: Ok_Toan11_TracNghiem_Lan2
\begin{ex}%[1D6H2-3]%[ma:88d4c9cc]
	Với $a$, $b$ là các số thực dương tùy ý và $a \ne 1$,  $\log_{\tfrac{1}{a}}\dfrac{1}{b^3}$ bằng
	\choice 
		{\True  $3\log_a b$}
		{$\log_a b$} 
		{$-3\log_a b$} 
		{$\dfrac{1}{3}\log_a b$} 
	 \loigiai{
		Ta có $\log_{\tfrac{1}{a}} \dfrac{1}{b^3}=\log_{a^{-1}}  b^{-3}=3\log_a b$.
	 }
\end{ex}

% nguồn: Ok_Toan11_TracNghiem_Lan2
\begin{ex}%[1D6H2-3]%[ma:aa28020a]
	Cho các số thực âm $a$, $b$ thỏa mãn$\log_{\sqrt 3}|a|-\log_3b^2=2$. Khẳng  định nào sau đây đúng?
	 \choice
		{$\dfrac{a}{b^2}=-9$} 
		{\True  $\dfrac{a}{b}=3$}
		 {$\dfrac{a}{b^2}=-3$}
		 {$\dfrac{a}{b}=-3$}
	 \loigiai{
		Ta có \[\log_{\sqrt 3}|a|-\log_3b^2=2\Leftrightarrow  \log_3a^2-\log_3b^2=2\Leftrightarrow \left(\dfrac{a}{b}\right)^2=9.\]
		Theo giả thiết $a$, $b$ là các số thực âm nên  $\dfrac{a}{b}>0$.\\
		Khi  đó $\left(\dfrac{a}{b}\right)^2=9\Leftrightarrow \dfrac{a}{b}=3$.
	} 
\end{ex}

% nguồn: Ok_Toan11_TracNghiem_Lan2
\begin{ex}%[1D6H2-3]%[ma:5f340d4f]
	Với  mọi $a$, $b$ thỏa mãn $\log_{\sqrt 2}(2a)+\log_8(b^9)=3$. Khẳng định nào dưới đây đúng?
	\choice 
		 {$a+b^9=8$}
		{\True  $a^2b^3=2$}
		{$ab^3=4$} 
		{$ab^9=8$} 
	 \loigiai{
		Ta có \[\log_{\sqrt 2}(2a)+\log_8b^9=3\Rightarrow \log_2(2a)^2+\log_2b^3=3\Rightarrow 4a^2b^3=8\Leftrightarrow  a^2b^3=2.\]
	} 
\end{ex}

% nguồn: Ok_Toan11_TracNghiem_Lan2
\begin{ex}%[1D6H2-3]%[ma:77ebbb54]
	Với  $a$, $b$ là các số thực dương tùy ý thỏa mãn $\log_3a-2\log_9b=2$, mệnh đề nào dưới đây đúng?
	\choice 
		 {$a=9b^2$}
		 {\True $a=9b$}
		 {$a=6b$}
		 {$b=9a$}
	 \loigiai{
		Ta  có \[\log_3a-2\log_9b=2\Leftrightarrow \log_3a-\log_3b=2\Leftrightarrow \log_3\left(\dfrac{a}{b}\right)=2\Leftrightarrow a=9b.\]
	} 
\end{ex}

% nguồn: Ok_Toan11_TracNghiem_Lan2
\begin{ex}%[1D6H2-3]%[ma:f2508e89]
	Với mọi $a$,  $b$ thỏa mãn $\log_3a^2+\log_3b=5$, khẳng định nào sau đây đúng?
	\choice 
		 {$a^2b=9$}
		{\True $a^2b=243$} 
		 {$a^2+b=243$}
		 {$a^3+b=15$}
	\loigiai{ 
		Ta có \[\log_3a^2+\log_3b=5\Leftrightarrow \log_3(a^2\cdot b)=5\Leftrightarrow a^2\cdot b=3^5.\] 
		 Vậy $a^2\cdot b=243$.
	 }
\end{ex}

% nguồn: Ok_Toan11_TracNghiem_Lan2
\begin{ex}%[1D6H2-3]%[ma:d6fe37f9]
	Với mọi số thực dương $a$, $b$, $c$ thỏa mãn $\log a^2+\log b-\log  c=0$, khẳng định nào sau đây đúng?
	 \choice
		{$a^2+b-c=0$} 
		{\True  $a^2b-c=0$}
		{$a^2+b-c=1$} 
		 {$a^2b-c=1$}
	 \loigiai{
		Ta có \[\log a^2+\log b-\log  c=0\Leftrightarrow \log (a^2b)=\log c\Leftrightarrow a^2b=c\Leftrightarrow a^2b-c=0.\]
	} 
\end{ex}

% nguồn: Ok_Toan11_TracNghiem_Lan2
\begin{ex}%[1D6H2-3]%[ma:b753e7f1]
	Cho $a$, $b$ là hai số thực dương thỏa mãn $\log_3a^5-\log_9b^2=4$, khẳng định  nào dưới đây đúng?
	 \choice
		 {$a^5-b=27$}
		 {$\dfrac{a^5}{b}=27$}
		 {$a^5-b^2=81$}
		{\True  $\dfrac{a^5}{b}=81$}
	\loigiai{ 
		Ta có 
		 \allowdisplaybreaks
		\begin{eqnarray*}
			\log_3a^5-\log_9b^2=4&\Leftrightarrow&\log_3a^5-\log_{3^2}b^2=4\\ 
			&\Leftrightarrow&\log_3a^5-\log_3b=4\\ 
			 &\Leftrightarrow&\log_3\dfrac{a^5}{b}=4\\
			 &\Leftrightarrow&\dfrac{a^5}{b}=81.
		\end{eqnarray*}
	 }
\end{ex}

% nguồn: Ok_Toan11_TracNghiem_Lan2
\begin{ex}%[1D6H2-3]%[ma:b3311b6b]
	Cho $a$, $b$ là các số  dương thỏa mãn $4\log_3a+7\log_3b=2$. Khẳng định nào sau đây là đúng?
	 \choice
		 {$4a+7b=2$}
		{$a^4b^7=2$} 
		{\True $a^4b^7=9$} 
		 {$4a+7b=9$}
	\loigiai{ 
		Ta có  \[4\log_3a+7\log_3b=2\Leftrightarrow \log_3a^4+\log_3b^7=2\Leftrightarrow \log_3a^4b^7=2\Leftrightarrow a^4b^7=9.\]
	} 
\end{ex}

% nguồn: Ok_Toan11_TracNghiem_Lan2
\begin{ex}%[1D6H2-3]%[ma:80402143]
	Giả sử $a$, $b$ là các số thực dương tùy ý thỏa mãn $a^2b^3=4^4$. Mệnh đề  nào sau đây là đúng?
	 \choice
		 {$2\log_2a+3\log_2b=4$}
		{\True  $2\log_2a+3\log_2b=8$}
		{$2\log_2a+3\log_2b=32$} 
		{$2\log_2a+3\log_2b=16$} 
	 \loigiai{
		Ta  có
		\[a^2b^3=4^4\Leftrightarrow \log_2(a^2b^3)=\log_24^4 \Leftrightarrow \log_2a^2+\log_2b^3=\log_22^8 \Leftrightarrow 2\log_2a+3\log_2b=8.\] 
	 }
\end{ex}

% nguồn: Ok_Toan11_TracNghiem_Lan2
\begin{ex}%[1D6H2-3]%[ma:2734bd6a]
	 Với $a$, $b$ là các số thực dương tuỳ ý và $a \ne 1$, $\log_{\tfrac{1}{a}} b^{-3}$ bằng
	\choice 
		{$\log_a  b$}
		{$-3\log_a b$} 
		{$\dfrac{1}{3}\log_a  b$}
		{\True $3\log_a b$} 
	\loigiai{ 
		 Ta có $\log_{\tfrac{1}{a}} b^{-3}=\log_{a^{-1}} b^{-3}=-1\cdot(-3) \log_a b=3\log_a b$.
	 }
\end{ex}

% nguồn: Ok_Toan11_TracNghiem_Lan2
\begin{ex}%[1D6H2-3]%[ma:2d961395]
	Với mọi $a$, $b$  thỏa mãn $\log_2a^3+\log_2b=5$. Khẳng định nào sau đây đúng?
	\choice 
		{$a^3+b=32$} 
		 {$a^3+b=25$}
		{\True $a^3b=32$} 
		{$a^3b=25$} 
	\loigiai{ 
		Ta có $\log_2a^3+\log_2b=5\Leftrightarrow  \log_2(a^3b)=5\Leftrightarrow a^3b=32$.
	 }
\end{ex}

% nguồn: Ok_Toan11_TracNghiem_Lan2
\begin{ex}%[1D6H2-3]%[ma:879fd1fc]
	Với  mọi $a$, $b$ thỏa mãn $\log_33a^2+\log_3b^3=4$, khẳng định nào dưới đây đúng?
	 \choice
		 {$a^2+b^3+1=81$}
		{\True  $a^2b^3=27$}
		{$a^2+b^3=27$} 
		 {$a^2b^3=81$}
	\loigiai{ 
		Ta có 
		\allowdisplaybreaks 
		\begin{eqnarray*}
			 \log_33a^2+\log_3b^3=4&\Leftrightarrow&1+\log_3a^2+\log_3b^3=4\\
			 &\Leftrightarrow&\log_3a^2+\log_3b^3=3\\
			 &\Leftrightarrow&\log_3a^2b^3=3\\
			 &\Leftrightarrow&a^2b^3=27.
		\end{eqnarray*}
	 }
\end{ex}

% nguồn: Ok_TracNghiemLop11_Lan1
\begin{ex}%[1D6H2-3]%[ma:47c25ae9]
	Cho $a > 0$ và $a\ne 1$ và $b > 0$.  Rút gọn biểu thức $P=\sqrt{\log_a^2(ab)-\dfrac{2\log b}{\log a}-1}$
	 \choice
		{$P=\left|\log_a  b+1\right|$}
		 {$P=\left|\log_a b-1\right|$}
		 {\True $P=\left|\log_a b\right|$}
		 {$P=0$}
	 \loigiai{
		 Ta có:
		\begin{eqnarray*}
			P&=&\sqrt{\log_a^2(ab)-\dfrac{2\log b}{\log a}-1}=\sqrt{\left[\log_a (ab)\right]^2-2\log_a  b-1}\\
			&=&\sqrt{\left(1+\log_a b\right)^2-2\log_a b-1}=\sqrt{\log_a^2b}=\left|\log_a b\right|. 
		\end{eqnarray*}
	 }
\end{ex}

% nguồn: Ok_TracNghiemLop11_Lan1
\begin{ex}%[1D6H2-3]%[ma:032faa9c]
	Cho $\log_62=a$, $\log_65=b$.  Tính $ I=\log_35$ theo $ a$, $ b$.
	 \choice
		{$  I=\dfrac{b}{1+a}$}
		{\True $ I=\dfrac{b}{1-a}$} 
		{$ I=\dfrac{b}{a-1}$} 
		{$  I=\dfrac{b}{a}$}
	 \loigiai{
	Ta có $\log_35=\dfrac{\log_65}{\log_63}=\dfrac{\log_65}{\log_66-\log_62}=\dfrac{b}{1-a}$.} 
\end{ex}

% nguồn: Ok_TracNghiemLop11_Lan1
\begin{ex}%[1D6H2-3]%[ma:3f7d4f8c]
	Cho $a$, $b$ là các số thực dương ($ a\ne 1$) thỏa mãn  $\log_{a^3} \dfrac{a^5}{\sqrt[4]{b}}=2$. Tính giá trị của biểu thức $\log_a b$.
	\choice 
		{$4$} 
		 {$\dfrac{1}{4}$}
		 {$-\dfrac{1}{4}$}
		{\True $-4$} 
	\loigiai{ 
		 Ta có $2=\log_{a^3} \dfrac{a^5}{\sqrt[4]{b}}=\dfrac{1}{3}\log_{a} \dfrac{a^5}{{b}^{\frac{1}{4}}}=\dfrac{1}{3}\left( \log_a a^5 -\log_a b^{\frac{1}{4}}\right)=\dfrac{1}{3}\left(5-\dfrac{1}{4}\log_a b\right)$.\\
		Từ đó, suy  ra $6=5-\dfrac{1}{4}\log_a b\Rightarrow \dfrac{1}{4}\log_a b=-1\Rightarrow \log_a b =-4$.
	} 
\end{ex}

% nguồn: Ok_TracNghiemLop11_Lan1
\begin{ex}%[1D6H2-3]%[ma:8bd6267b]
	Với mọi số thực dương $a$ và $b$ thỏa mãn $a^2+b^2=8ab$, mệnh  đề nào dưới đây đúng?
	 \choice
		{$\log (a+b)=\dfrac{1}{2}  (\log a+\log b)$}
		 {$\log (a+b)=\dfrac{1}{2}+\log a+\log b$}
		{\True $\log (a+b)=\dfrac{1}{2} (1+\log  a+\log b)$}
		{$\log (a+b)=1+\log  a+\log b$}
	\loigiai{ 
		 Ta có $a^2+b^2=8ab\Leftrightarrow (a+b)^2=10ab$.\\
		Lấy log cơ  số $10$ hai vế ta được: $\log (a+b)^2=\log (10ab)\Leftrightarrow 2\log (a+b)=\log 10+\log a+\log b$.\\
		Hay $\log (a+b)=\dfrac{1}{2} (1+\log a+\log b)$. 
	} 
\end{ex}

% nguồn: Ok_TracNghiemLop11_Lan1
\begin{ex}%[1D6H2-3]%[ma:5adeb7d3]
	Cho các số thực $ a$, $ b$ thỏa  $ 0 < a\ne 1$ và $ b\ne 0$. Rút gọn biểu thức $$ A=\log_a{b^2}+\log_{a^2}{b^4}+\log_{a^6}{b^6}.$$
	\choice 
		{$  A=5\log_ab$}
		 {$ A=5\log_ba$}
		{\True $  A=5\log_a\left| b\right|$}
		{$ A=5$} 
	 \loigiai{
	 Ta có $ A=2\log_a\left| b\right|+2\log_a\left| b\right|+\log_a\left| b\right|=5\log_a\left| b\right|$.}
\end{ex}

% nguồn: Ok_TracNghiemLop11_Lan1
\begin{ex}%[1D6H2-3]%[ma:c140591d]
	Cho $a$ là  số thực dương thỏa mãn $\log_3a=5$. Giá trị của biểu thức $P=\log_3(9a)$ bằng
	 \choice
		{$10$} 
		 {$45$}
		 {\True $7$}
		 {$3$}
	\loigiai{ 
		Ta  có $P=\log_3(9a)=2+\log_3a=2+5=7$.
	} 
\end{ex}

% nguồn: Ok_TracNghiemLop11_Lan1
\begin{ex}%[1D6H2-3]%[ma:1f58f320]
	Cho $\log_8 |x|+\log_4 y^2=5$ và $\log_8 |y|+\log_4 x^2=7$. Tìm giá trị  của biểu thức $P=|x|-|y|$.
	\choice 
		{\True  $P=56$}
		{$P=16$} 
		{$P=8$} 
		 {$P=64$}
	\loigiai{ 
		 Điều kiên: $x,y\ne 0$.\\
		Cộng vế  với vế của hai phương trình, ta được:
		\[\log_8 |xy|+\log_4 x^2y^2=12\Leftrightarrow \log_2  |xy|=9\Leftrightarrow |xy|=512.\qquad (1)\]
		Trừ  vế với vế của hai phương trình, ta được:
		\[\log_8 |\dfrac{x}{y} |+\log_4 \dfrac{y^2}{x^2}=-2\Leftrightarrow \log_2 |\dfrac{x}{y} |=3\Leftrightarrow  |\dfrac{x}{y} |=8\Leftrightarrow |x|=8|y|. \qquad (2)\]
		 Từ (1) và (2) suy ra $|y|=8\Rightarrow |x|=64\Leftrightarrow P=56$.
	} 
\end{ex}

% nguồn: Ok_TracNghiemLop11_Lan1
\begin{ex}%[1D6H2-3]%[ma:4f29e158]
	Cho các số thực dương $a$, $b$ với $a\ne 1$. Khẳng định nào sau đây là  khẳng định đúng?
	 \choice
		 {$\log _{a^2} (ab)=\dfrac{1}{2} \log_a b$}
		{\True $\log _{a^2} (ab)=\dfrac{1}{2}+\dfrac{1}{2} \log_a  b$}
		 {$\log _{a^2} (ab)=2+\log_a b$}
		{$\log _{a^2} (ab)=\dfrac{1}{4}  \log_a b$}
	\loigiai{ 
		 Ta có $\log _{a^2} (ab)=\dfrac{1}{2} \left(\log_a (ab)\right)=\dfrac{1}{2} \left(\log_a a+\log_a b\right)
		 =\dfrac{1}{2}+\dfrac{1}{2} \log_a b$.
	} 
\end{ex}

% nguồn: Ok_TracNghiemLop11_Lan1
\begin{ex}%[1D6H2-3]%[ma:95110287]
	 Cho $a < 0, b > 0$. Mệnh đề nào sau đây \textbf{sai}?
	 \choice
		{$\sqrt[4]{a^4 b^8}=a b^2$} 
		{$\sqrt[4]{a^4 b^8}=\left|a b^2\right|$} 
		{$\sqrt[4]{a^4 b^8}=|a|b^2$} 
		{\True $\sqrt[4]{a^4  b^8}=-a b^2$}
	\loigiai{ 
		Ta có $\sqrt[4]{a^4 b^8}=\left|a b^2\right| = |a|\cdot \left|b^2\right| =  -ab^2$.
	 }
\end{ex}

% nguồn: Ok_TracNghiemLop11_Lan1
\begin{ex}%[1D6H2-3]%[ma:48a888d0]
	Cho $a$, $b$ là các  số thực dương ($ a\ne 1$) thỏa mãn $\log_2{a}=2$, $ \log_4 b=3$. Tính giá trị của biểu thức $P=\log_a (a^2b)$.
	 \choice
		{$P=10$} 
		{\True  $P=5$}
		{$P=2$} 
		 {$P=1$}
	 \loigiai{
		Ta có $\heva{& \log_2 a=2\\& \log_4 b=3}\Leftrightarrow  \heva{&a=2^2\\&b=4^3}\Leftrightarrow \heva{&a=4\\&b=4^3.}$\\
		 Vậy $P=\log_a(a^2b)=\log_4(4^2\cdot 4^3)=\log_4 4^5=5$.
	 }
\end{ex}

% nguồn: Ok_TracNghiemLop11_Lan1
\begin{ex}%[1D6H2-3]%[ma:5bc9d4b5]
	Với $a$ là số thực dương tùy ý, $\log _2a^2+\log  _4a$ bằng
	\choice 
		{$\dfrac{3}{2}\log  _2a$}
		{\True $\dfrac{5}{2}\log  _2a$}
		{$\log  _2a$}
		 {$\dfrac{1}{2}\log _2a$}
	 \loigiai{
	 Ta có $\log _2a^2+\log _4a=2\log _2a+\dfrac{1}{2}\log _2a=\dfrac{5}{2}\log _2a$.}
\end{ex}

% nguồn: Ok_TracNghiemLop11_Lan1
\begin{ex}%[1D6H2-3]%[ma:2e590e1d]
	 Cho $a$, $b$ là các số thực dương thỏa mãn $\log _{\sqrt{ab}} \left(a\sqrt[3]b \right)=3$. Tính $\log _{\sqrt{ab}} \left(b\sqrt[3]a \right)$.
	\choice 
		 {$-3$}
		 {$\dfrac{1}{3}$}
		 {$3$}
		{\True  $-\dfrac{1}{3}$}
	\loigiai{ 
		 Ta có
		\begin{align*}
			\log _{\sqrt{ab}} \left(a\sqrt[3]b \right)=3&\Leftrightarrow \dfrac{\log_a \left(a\sqrt[3]b  \right)}{\log_a \sqrt{ab}}=3\\
			&\Leftrightarrow \dfrac{\log_a a+\dfrac{1}{3}  \log_a b}{\dfrac{1}{2} \left(\log_a a+\log_a b\right)}=3\\
			 &\Leftrightarrow \dfrac{2}{3}\cdot \dfrac{3+\log_a b}{1+\log_a b}=3\\
			&\Leftrightarrow 6+2\log_a  b=9+9\log_a b\\
			&\Leftrightarrow  \log_a b=-\dfrac{3}{7}.
		\end{align*}
		Khi đó 
		$\log _{\sqrt{ab}} \left(b\sqrt[3]a \right)=\dfrac{\log_a \left(b\sqrt[3]a  \right)}{\log_a \sqrt{ab}}=\dfrac{\log_a b+\dfrac{1}{3} \log_a a}{\dfrac{1}{2} \left(\log_a a+\log_a b\right)}=\dfrac{-\dfrac{3}{7}+\dfrac{1}{3}}{\dfrac{1}{2} \left(1-\dfrac{3}{7} \right)}=-\dfrac{1}{3}$.
	} 
\end{ex}

% nguồn: Ok_TracNghiemLop11_Lan1
\begin{ex}%[1D6H2-3]%[ma:0d4a2814]
	Xét tất cả các  số dương $a$ và $b$ thỏa mãn $\log_2 a=\log_8 (ab)$. Mệnh đề nào dưới đây đúng?
	 \choice
		 {$a=b^2$}
		{$a^3=b$} 
		{$a=b$} 
		{\True $a^2=b$} 
	\loigiai{ 
		Theo đề  ta có:
		\begin{eqnarray*}
			 && \log_2a=\log_8(ab)\Leftrightarrow \log_2a=\dfrac{1}{3} \log_2(ab)\Leftrightarrow 3\log_2a=\log_2(ab) \\
			&\Leftrightarrow& \log_2a^3=\log_2(ab)\Leftrightarrow a^3=ab\Leftrightarrow a^2=b. 
		\end{eqnarray*}
	} 
\end{ex}

% nguồn: Ok_TracNghiemLop11_Lan1
\begin{ex}%[1D6H2-3]%[ma:70e65614]
	Cho  $a$, $b$, $c> 1$. Biết rằng biểu thức $P=\log_a (bc)+\log_b (ac)+4\log_c (ab)$ đạt giá trị lớn nhất $m$ khi $\log_b c=n$. Tính giá trị $m+n$.
	\choice 
		{$m+n=10$} 
		 {$m+n=\dfrac{25}{2}$}
		 {$m+n=14$}
		{\True $m+n=12$} 
	\loigiai{ 
		Ta có $\begin{aligned}
		P &=  \log_a b+\log_a c+\log_b a+\log_b c+4\log_c a+4\log_c b \Leftrightarrow P\\
		& = \left(\log_a b+\dfrac{1}{\log_a b}  \right)+\left(\log_a c+\dfrac{4}{\log_a c} \right)+\left(\log_b c+\dfrac{4}{\log_b c} \right) \ge 2+4+4\\
		& =  10.
	\end{aligned}$\\
	Suy ra  $m= 10$.\\
	Dấu đẳng xảy ra khi  $\log_a b=1$, $\log_a c=2$, $\log_b c=2 \Rightarrow n=2$.\\
	 Vậy $m+n=12$.
 }
\end{ex}

% nguồn: Ok_TracNghiemLop11_Lan1
\begin{ex}%[1D6H2-3]%[ma:02a33789]
	Cho $a$, $b$ là các số thực dương ($ a\ne 1$) thỏa mãn $\log_a b=2$. Tính giá trị  của biểu thức
	$P=\log_{\frac{\sqrt{a}}{b}}\left(  a\sqrt[3]{b}\right)$.
	 \choice
		 {$P=\dfrac{2}{15}$}
		{$P=-\dfrac{2}{9}$} 
		{\True  $P=-\dfrac{10}{9}$}
		{$P=\dfrac{2}{3}$} 
	 \loigiai{
		Ta  có
	$$P=\log_{\frac{\sqrt{a}}{b}}\left( a\sqrt[3]{b}\right)= \dfrac{\log_a\left(a\sqrt[3]{b}\right)}{\log_a\left(\dfrac{\sqrt{a}}{b}\right)}=\dfrac{\log_a a+\log_a b^{\frac{1}{3}}}{\log_a a^{\frac{1}{2}-\log_a b}}=\dfrac{1+\dfrac{1}{3}\log_a b}{\dfrac{1}{2}-\log_a b}=\dfrac{1+\dfrac{2}{3}}{\dfrac{1}{2}-2}=-\dfrac{10}{9}.$$} 
\end{ex}

% nguồn: Ok_TracNghiemLop11_Lan1
\begin{ex}%[1D6H2-3]%[ma:68849b4c]
	Cho các số thực  dương $a$, $b$, $x$ thỏa mãn $\log_3 x=4\log_3 a+7\log_3 b$. Mệnh đề nào dưới đây đúng
	\choice{$x=4a+7b$}{$x=4a-7b$}{\True $x=a^4b^7$}{$x=a^{\tfrac{1}{4}}b^{\tfrac{1}{7}}$} 
	\loigiai{ 
		Ta có \begin{align*}
		 &\ \log_3 x=4\log_3 a+7\log_3 b\\
		 \Leftrightarrow&\ x=3^{4\log_3 a+7\log_3 b}\\
		\Leftrightarrow&\ x=3^{\log_3  a^4+\log_3 b^7}\\
		\Leftrightarrow&\ x= 3^{\log_3 a^4}\cdot 3^{\log_3  b^7}\\
		 \Leftrightarrow&\ x=a^4 b^7.
	\end{align*}
} 
\end{ex}

% nguồn: Ok_TracNghiemLop11_Lan1
\begin{ex}%[1D6H2-3]%[ma:ceee9193]
	Cho các số thực dương $ a,b$ và $ a\ne 1$. Biểu thức $\log_a{a^2}b$  bằng
	 \choice
		 {\True $ 2+\log_ab$}
		{$ 1+\log_ab$} 
		{$ 2\left(1+\log_ab\right)$} 
		{$ 2\log_ab$} 
	\loigiai{ 
	Ta có $\log_a{a^2}b=\log_a{a^2}+\log_ab=2\log_aa+\log_ab=2+\log_ab$.} 
\end{ex}

% nguồn: Ok_TracNghiemLop11_Lan1
\begin{ex}%[1D6H2-3]%[ma:89e1c6aa]
	Cho $2$ số thực dương $a$, $b$ thỏa mãn $\sqrt{a} \ne b$, $a\ne 1$, $\log_a b=2$. Tính $T=\log _{\tfrac{\sqrt{a}}{b}} \sqrt[3]{ba}$ 
	 \choice
		{$T=-\dfrac{2}{5}$} 
		{$T=\dfrac{2}{5}$} 
		 {$T=\dfrac{2}{3}$}
		 {\True $T=-\dfrac{2}{3}$}
	\loigiai{ 
		Ta có: $\log_a b=2\Rightarrow \log_b a=\dfrac{1}{2}$. 
		\begin{align*}
			&T=\log  _{\tfrac{\sqrt{a}}{b}} \sqrt[3]{ba}=\log _{\tfrac{\sqrt{a}}{b}} \sqrt[3]b+\log _{\tfrac{\sqrt{a}}{b}} \sqrt[3]a\\
			&=\dfrac{1}{\log _{\sqrt[3]b} \tfrac{\sqrt{a}}{b}}+\dfrac{1}{\log _{\sqrt[3]a}  \tfrac{\sqrt{a}}{b}}\\
			&=\dfrac{1}{\log _{\sqrt[3]b} \sqrt{a}-\log _{\sqrt[3]b} b}+\dfrac{1}{\log _{\sqrt[3]a}  \sqrt{a}-\log _{\sqrt[3]a} b}\\
			&=\dfrac{1}{\dfrac{3}{2} \log_b  a-3}+\dfrac{1}{\dfrac{3}{2}-3\log_a b}\\
			&=\dfrac{1}{\dfrac{3}{2}\cdot \dfrac{1}{2}-3}+\dfrac{1}{\dfrac{3}{2}-3\cdot 2}=-\dfrac{2}{3}. 
		\end{align*}
	} 
\end{ex}

% nguồn: Ok_TracNghiemLop11_Lan1
\begin{ex}%[1D6H2-3]%[ma:ece3549e]
	Cho $0<a\ne 1$. Mệnh đề nào  sau đây {\bf sai}?
	\choice 
		{$\log_a a=1$} 
		{\True $a^{\log_5  a}=a$}
		{$5^{\log_5 a}=a$} 
		{$\log_a  1=0$}
	 \loigiai{
		$a^{\log_5 a}=a$ là khẳng  định sai.
	} 
\end{ex}

% nguồn: Ok_TracNghiemLop11_Lan1
\begin{ex}%[1D6H2-3]%[ma:0dd65927]
	Nếu $\log_2x=5\log_2a+4\log_2b$ ($ a,b  > 0$) thì $ x$ bằng
	 \choice
		{$  4a+5b$}
		 {\True $a^5b^4$}
		{$a^4b^5$} 
		 {$ 5a+4b$}
	\loigiai{ 
	Ta có $\log_2x=5\log_2a+4\log_2b\Leftrightarrow\log_2x=\log_2a^5b^4\Leftrightarrow x=a^5b^4$.} 
\end{ex}

% nguồn: Ok_TracNghiemLop11_Lan1
\begin{ex}%[1D6H2-3]%[ma:9723a04d]
	Cho $\log_3  2=a$; $\log_3 5=b$ , khi đó $\log_340$ bằng
	 \choice
		{$3\rm{a}-b$} 
		{$a-3b$} 
		 {\True $3\rm{a}+b$}
		 {$a+3b$}
	 \loigiai{
	Ta có $\log_  340=\log_3\left(2^3\cdot 5\right)=3\log_3 2+\log_3 5=3a+b$.}
\end{ex}

% nguồn: Ok_TracNghiemLop11_Lan1
\begin{ex}%[1D6H2-3]%[ma:c13af606]
	Với $a$ là số thực dương tùy ý, biểu thức $T=\ln  a^2-\ln\sqrt[3]{a}$ bằng
	 \choice
		{$\dfrac{4}{3}\ln  a$}
		 {$\ln\dfrac{3}{5}$}
		{$\dfrac{5}{3}$} 
		{\True  $\dfrac{5}{3}\ln a$}
	 \loigiai{
		Ta có $T=\ln a^2-\ln\sqrt[3]{a} =  2\ln a - \dfrac{1}{3}\ln a = \dfrac{5}{3}\ln a$.
	} 
\end{ex}

% nguồn: Ok_TracNghiemLop11_Lan1
\begin{ex}%[1D6V2-3]%[ma:64ec2ded]
	Cho $a$ là hai số thực dương khác $1$. Đặt  $\log_3a=m$. Tính theo $m$ giá trị của biểu thức $D=\log_{\tfrac{1}{3}} a-\log_{\sqrt{3}} a+\log_a 9$.
	\choice 
		{\True $D=\dfrac{2-3m^2}{m}$} 
		 {$D=\dfrac{3m^2-2}{m}$}
		 {$D=\dfrac{4-3m^2}{2m}$}
		 {$D=-3m$}
	\loigiai{Ta có $\begin{aligned}
	D&\ =\ \log_{\tfrac{1}{3}} a-\log_{\sqrt{3}} a+\log_a  9\\
	&\ =\ \log_{3^{-1}} a- \log_{3^{\frac{1}{2}}} a+ \log_a 3^2\\ 
	&\ =\ -\log_3 a-2 \log_3 a+2 \log_a  3\\
	&\ =\ -3 \log_3 a+\dfrac{2}{\log_3 a}\\ 
	&\  =\ -3 m+\dfrac{2}{m}\\
	&\ =\  \dfrac{2-3 m^2}{m}.
\end{aligned}$
 }
\end{ex}

% nguồn: Ok_TracNghiemLop11_Lan1
\begin{ex}%[1D6V2-3]%[ma:784e25e1]
	Cho $a=\log_{30}  3$ và $b=\log_{30} 5$. Tính $\log_{30} 1350$ theo $a$ và $b$.
	 \choice
		 {$1+2a-b$}
		 {\True $1+2a+b$}
		{$1-2a+b$} 
		 {$-1+2a+b$}
	\loigiai{Ta  có $\log_{30} 1350=\log_{30}( 30\cdot 3^2\cdot 5)=1+2\log_{30} 3+\log_{30} 5=1+2a+b$.
	 }
\end{ex}

% nguồn: Ok_TracNghiemLop11_Lan1
\begin{ex}%[1D6V2-3]%[ma:d04fa7f7]
	Cho $\log_25=a$, $\log_35=b$.  Hãy biểu diễn $\log_65$ theo $a$ và $b$.
	 \choice
		{$\log_65=\dfrac{1}{a+b}$} 
		{\True $\log_65=\dfrac{a  b}{a+b}$}
		{$\log_65=a+b$} 
		{$\log_65=a^2+b^2$} 
	\loigiai{Ta có $\log_65=\dfrac{1}{\log_56}=\dfrac{1}{\log_5(2.3)}=\dfrac{1}{\log_52+\log_53}=\dfrac{1}{\dfrac{1}{a}+\dfrac{1}{b}}=\dfrac{a b}{a+b}$. 
	} 
\end{ex}


%%%%%%%%%%%%%%%% ĐÚNG/SAI %%%%%%%%%%%%%%%%

% nguồn: Ok_dungsai_lop11_Lan1
\begin{ex}%[1D6H2-3]%[ma:0020d607]
 Cho các biểu thức $P=\dfrac{\log_{a}(a^3b^2)-\log_b\left(\dfrac{b^3}{a^2}\right)}{\log_a^2 b+1}$ và $Q=\log_a b^3+\log_{a^2}b^6$  với $a$, $b$ là các số dương và $a\ne 1$.
 \choiceTF 
   {\True $Q=6\log_{a}b$}
  {$P= 6 \log_b a$} 
   {$Q=3P$}
  {\True $Q\cdot P=12$} 
  \loigiai{ Ta có
  \begin{eqnarray*}
   P&=& \dfrac{\log_a a^3  +\log_{a} b^2-\left( \log_b b^3 -\log_{b} a^2\right)}{\log_{a}^2b +1}\\
   &=&\dfrac{3+2\log_{a} b -3  +2\log_b a}{\log_a^2 b +1}\\
   &=& \dfrac{2\left( \log_a b +\dfrac{1}{\log_a b}\right)}{\log_{a}^2b +1}\\ 
   &=&\dfrac{2\left( \dfrac{\log_a^2 b+1}{\log_a  b }\right)}{\log_a^2b +1}
    = \dfrac{2}{\log_a b}=2\log_b a.
  \end{eqnarray*}
  $Q=\log_a b^3+\log_{a^2}b^6=3\log_{a} b+\dfrac{6}{2}\cdot \log_{a} b =6\log_{a} b$. 
  \begin{itemchoice}
   \itemch  Đúng. $Q=6\log_{a} b$.
    \itemch Sai. $P=2\log_b a$.
    \itemch Sai. $Q=\dfrac{3}{P}$.
   \itemch Đúng. $QP= \log_{a} b\cdot 2\log_b  a =12$.
  \end{itemchoice}
  }
\end{ex}

% nguồn: Ok_dungsai_lop11_Lan1
\begin{ex}%[1D6H2-3]%[5-KNTT]%[ma:3b155bf7]
 Cho biểu thức $ Q=2^{\log_{16}{x^4}+\log_2x^2}$ với $  x$ là số thực khác $0$. Vậy
 \choiceTF 
  {\True $ Q > 0$} 
  {\True Khi  $ x=2$ thì $ Q=8$}
  {Khi $ x=-2$  thì $ Q=-8$}
  {Khi $ x=3$  thì $ Q=9$}
 \loigiai{ 
  Ta  có $Q=2^{\log_{16}{x^4}+\log_2x^2}=2^{\log_2|x|+\log_2x^2}=2^{\log_2|x|x^2}=|x|x^2$.
  \begin{itemize}
   \item $Q=|x|x^2>0$ với mọi  $x\ne 0$.
   \item Với $x= 2$ thì  $Q=|2|\cdot 2^2=8$.
   \item Với  $x= -2$ thì $Q=|-2|\cdot (-2)^2=8$.
   \item  Với $x=3$ thì $Q=|3|\cdot 3^2=27$.
  \end{itemize}
  }
\end{ex}

% nguồn: Ok_dungsai_lop11_Lan1
\begin{ex}%[1D6H2-3]%[ma:669ae7b6]
 Cho biểu thức $Q=2^{\log_{16}{x^4}+\log_2{x^2}}$ với $x$ là  số thực khác $0$.
 \choiceTF 
   {$Q=x^3$}
   {$Q>0$}
  {Khi $x=2$ thì $Q=8$} 
   {Khi $x=-2$ thì $\log_2 Q=-3$}
 \loigiai{ 
  Ta  có $Q=2^{\log_{16}{x^4}+\log_2{x^2}}=2^{4\cdot \tfrac{1}{4}\log_2|x|+2\log_2|x|}=2^{3\log_2|x|}=\left(2^{\log_2|x|}\right)^3=|x|^3$.
  \begin{itemchoice}
   \itemch Sai. Ta có $Q=|x|^3$  với mọi $x\ne 0$.
    \itemch Đúng. $Q=|x|^3>0$, với mọi $x\ne 0$.
   \itemch Đúng. Khi $x=2$ thì  $Q=8$.
   \itemch Sai. Khi $x=-2$ thì  $\log_2 Q=\log_2 |-2|^3=3$
  \end{itemchoice}
  }
\end{ex}

% nguồn: Ok_dungsai_lop11_Lan1
\begin{ex}%[1D6H2-3]%[ma:4d782820]
 Cho biểu thức $A=\left(a^3\sqrt{a}\right)^{\log_a b}+\left(\sqrt[3]{b^2}\right)^{\log_b a}$ với $a, b >0$, $a\ne 1$, $b\ne 1$ và $B=\log \dfrac{a}{b}+\log \dfrac{b}{c}+\log \dfrac{c}{d}-\log \dfrac{a}{d}$  với $a$, $b$, $c$, $d$ là các số dương.
  \choiceTF
  {$A=\sqrt[3]{a}+\sqrt{b^4}$} 
  {$B=\dfrac{a}{b}$} 
  {\True  $A+B\sqrt{a}=\sqrt[3]{a^2}+\sqrt{b^7}$}
   {$A\sqrt[3]{a}-B\sqrt{b}=2a+\sqrt{b^7}$}
  \loigiai{
  \begin{itemchoice}
   \itemch Sai. Ta  có
   \begin{eqnarray*}
    A&=&\left(a^3\cdot  a^{\frac{1}{2}}\right)^{\log_a b}+ \left( b^{\frac{2}{3}}\right)^{\log_b a} \\
    &=&\left( a^{\frac{7}{2}}\right)^{\log_a b}+\left( b^{\frac{2}{3}}\right)^{\log_b  a}\\
    &=&\left(a^{\log_a b}\right)^{\frac{7}{2}} +\left(  b^{\log_b{a}}\right)^{\frac{2}{3}}\\
    &=&b^{\frac{7}{2}}+a^{\frac{2}{3}}=\sqrt[3]{a^2}+\sqrt{b^7}. 
   \end{eqnarray*}
   \itemch Sai.  $B=\log \left(\dfrac{a}{b}\cdot \dfrac{b}{c}\cdot \dfrac{c}{d}\right)-\log\dfrac{a}{d}=\log \left(\dfrac{a}{d}:\dfrac{a}{d}\right)=\log 1=0$.
    \itemch Đúng. $A+B\sqrt{a}=\sqrt[3]{a^2}+\sqrt{b^7}+0\sqrt{a}=\sqrt[3]{a^2}+\sqrt{b^7}$.
    \itemch Sai. $A\sqrt[3]{a}-B\sqrt{b}=\sqrt[3]{a}\cdot\sqrt[3]{a^2}+\sqrt[3]{a}\cdot\sqrt{b^7}-0\sqrt{b}=a+\sqrt[3]{a}\cdot\sqrt{b^7}$.
 \end{itemchoice} }
\end{ex}

% nguồn: Ok_dungsai_lop11_Lan1
\begin{ex}%[1D6H2-3]%[ma:bbad34dd]
 Cho các biểu thức sau: $A=\left(a^3 \sqrt{a}\right)^{\log _a b}+\left(\sqrt[3]{b^2}\right)^{\log _b a}$ với $\heva{&a, b>0 \\ &a \neq 1,  b \neq 1}$ và\\
 $B=\log \dfrac{a}{b}+\log \dfrac{b}{c}+\log \dfrac{c}{d}-\log \dfrac{a}{d}$ với $a$, $b$,  $c$, $d$ là các số dương. Khi đó
 \choiceTF 
   {$A=\sqrt[3]{a}+\sqrt{b^4}$}
  {$B=\dfrac{a}{b}$} 
  {\True $A+B  \sqrt{a}=\sqrt[3]{a^2}+\sqrt{b^7}$}
  {$A-B  \sqrt{b}=2 \sqrt[3]{a^2}+\sqrt{b^7}$}
  \loigiai{
  \begin{itemize}
   \item  Ta có
   \begin{eqnarray*}
    A&=& \left(a^3 \cdot a^{\tfrac{1}{2}}\right)^{\log _a b}+\left(b^{\tfrac{2}{3}}\right)^{\log  _b a}\\
    &=& \left(a^{\tfrac{7}{2}}\right)^{\log _a b}+\left(b^{\tfrac{2}{3}}\right)^{\log  _b a}\\
    &=& \left(a^{\log _a b}\right)^{\tfrac{7}{2}}+\left(b^{\log  _b a}\right)^{\tfrac{2}{3}}\\
     &=&b^{\tfrac{7}{2}}+a^{\tfrac{2}{3}}\\
     &=&\sqrt[3]{a^2}+\sqrt{b^7}.
   \end{eqnarray*}
   \item Ta có $B=\log \left(\dfrac{a}{b} \cdot \dfrac{b}{c} \cdot \dfrac{c}{d}\right)-\log  \dfrac{a}{d}=\log \left(\dfrac{a}{d}: \dfrac{a}{d}\right)=\log 1=0$.
   \item $A+B  \sqrt{a}=\sqrt[3]{a^2}+\sqrt{b^7}$.
   \item  $A-B \sqrt{b}=\sqrt[3]{a^2}+\sqrt{b^7}$.
  \end{itemize}
 } 
\end{ex}

% nguồn: Ok_dungsai_lop11_Lan1
\begin{ex}%[1D6H2-3]%[ma:d77127de]
 Cho  biểu thức $A=\log_2 x^2 +\log_{\frac{1}{2}}x^3+\log_4 x$, với $x$ là số thực dương.
 \choiceTF 
   {Rút gọn $A= \dfrac{1}{2}\log_2 x$}
   {\True Khi $\log_2 x =1$ thì $A=-\dfrac{1}{2}$}
  {\True Khi $x =4$ thì $\ln A$ không  tồn tại}
  {Khi $\log_2 x =-\sqrt{2}$ thì $\log_2 A=\dfrac{\sqrt{2}}{2}$} 
  \loigiai{
  \begin{eqnarray*}
   A&=& \log_2 x^2 +\log_{\frac{1}{2}}x^3+\log_4  x\\
    &=& 2\log_2 x +3\log_{2^{-1}} x+\log_{2^2}x\\
   &=&  2\log_2 x =3\log_2 x+\dfrac{1}{2}\log_2 x\\
   &=& -\dfrac{1}{2}\log_2 x. 
  \end{eqnarray*}
  \begin{itemchoice}
   \itemch Sai. Ta có  $A= -\dfrac{1}{2}\log_2 x$.
   \itemch Đúng. Khi $\log_2 x =1$ thì  $A=-\dfrac{1}{2}$.
   \itemch Đúng. Khi $x =4$ thì $\log_2 x=2$ và $A=-1$. Khi đó $\ln A$ không tồn  tại.
   \itemch Sai. Khi $\log_2 x =-\sqrt{2}$ thì $A=\dfrac{\sqrt{2}}{2}=\dfrac{1}{\sqrt{2}}=2^{-\tfrac{1}{2}}$ và  $\log_2 A =\log_2\left(2^{-\tfrac{1}{2}}\right)=-\dfrac{1}{2}$.
  \end{itemchoice}
  }
\end{ex}

% nguồn: Ok_dungsai_lop11_Lan2
\begin{ex}%[1D6H2-3]%[ma:ed340d6b]
	Cho  $a$, $b$, $c$ là các số thực dương tuỳ ý và $a\ne 1$.
	\choiceTF 
		{\True $\log_a  \left(a\cdot b\right)=1+\log_a b$}
		{$\log_a  \left(\dfrac{a^3}{b^2} \right)=\dfrac{3}{2\log_a b}$}
		{$\log_a  \left(b\cdot c\right)=\log_a b\cdot \log_a c$}
		{$\log_a  b+2\log_a c-\log_a 2=\log_a \left(b+c^2-2\right)$}
	\loigiai{ 
		\begin{itemchoice}
			 \itemch
			Ta có $\log_a\left(a\cdot b\right)=\log_a a+\log_a b=1+\log_a b$.  Vậy $\log_a\left(a\cdot b\right)=1+\log_a b$.
			 \itemch
			Ta có $\log_a \left(\dfrac{a^3}{b^2} \right)=\log_a a^3-\log_a b^2=3-2\log_a b$. Vậy $\log_a \left(\dfrac{a^3}{b^2}  \right)=\dfrac{3}{2\log_a b}$.
			\itemch 
			Ta có $\log_a\left(b\cdot c\right)=\log_a b+\log_a c$. Vậy $\log_a\left(b\cdot c\right)=\log_a b\cdot\log_a c$ 
			 \itemch
			Ta có $\log_a b+2\log_a c-\log_a 2=\log_a b+\log_a c^2-\log_a  2=\log_a \left(\dfrac{b\cdot c^2}{2} \right)$.\\
			Vậy $\log_a b+2\log_a c-\log_a 2=\log_a  \left(\dfrac{b\cdot c^2}{2} \right)$.
		\end{itemchoice}
	} 
\end{ex}

% nguồn: Ok_dungsai_lop11_Lan1
\begin{ex}%[1D6V2-3]%[ma:26a3949f]
 Cho các biểu thức sau: $A = \left( a^3 \sqrt{a} \right)^{\log_{a}{b}} + \left( \sqrt[3]{b^2} \right)^{\log_{b}{a}}$ với $\heva{&a,b>0\\ &a\ne 1, \, b\ne 1}$  \\
 và $B = \log\dfrac{a}{b} + \log\dfrac{b}{c} +  \log\dfrac{c}{d} - \log\dfrac{a}{d}$ với $a$, $b$, $c$, $d$ là các số dương. Khi đó
 \choiceTF 
   {$A=\sqrt[3]{a} + \sqrt{b^4}$}
   {$B=\dfrac{a}{b}$}
  {\True  $A+B\sqrt{a} = \sqrt[3]{a^2} + \sqrt{b^7}$}
  {$A-B\sqrt{b} = 2\sqrt[3]{a^2} +  \sqrt{b^7}$}
 \loigiai{ 
  Ta  có
  \allowdisplaybreaks 
  \begin{eqnarray*}
   A &=& \left( a^3\cdot a^{\frac{1}{2}}  \right)^{\log_{a}{b}} + \left( b^{\frac{2}{3}} \right)^{\log_{b}{a}} \\
   &=& \left( a^{\frac{7}{2}} \right)^{\log_{a}{b}} + \left( b^{\frac{2}{3}} \right)^{\log_{b}{a}}\\ &=&  \left( a^{\log_{a}{b}} \right)^{\frac{7}{2}} + \left( b^{\log_{b}{a}} \right)^{\frac{2}{3}} \\
   &=& b^{\frac{7}{2}} + a^{\frac{2}{3}}  \\
    &=& \sqrt[3]{a^2} + \sqrt{b^7}.
  \end{eqnarray*}
  $$ B = \log\left( \dfrac{a}{b}\cdot \dfrac{b}{c}\cdot \dfrac{c}{d} \right) - \log\dfrac{a}{d} = \log\left( \dfrac{a}{d} : \dfrac{a}{d} \right) = \log 1  = 0.$$
 } 
\end{ex}

% nguồn: Ok_dungsai_lop11_Lan1
\begin{ex}%[1D6V2-3]%[ma:af426971]
 Cho các biểu thức sau: ${P=\dfrac{\log_a\left(a^3 b^2\right)-\log_b\left(\dfrac{b^3}{a^2}\right)}{\log_a^2 b+1}}$ và $Q=\log_ab^3+\log_{a^2}b^6$ với $a$, $b$ là các  số dương và $a$ khác $1$. Vậy
  \choiceTF
  {\True  $Q=6\log_ab$}
   {$P=6\log_ba$}
   {$Q=3P$}
  {\True $QP=12$} 
  \loigiai{
  \begin{itemize}
   \item  Ta có $Q=3\log_ab+6\cdot \dfrac{1}{2}\log_ab=6\log_ab$.
   \item Ta  có
   \allowdisplaybreaks{\begin{eqnarray*}
   P & =&\dfrac{\log_a a^3+\log_a b^2-\left(\log_b b^3-\log_b a^2\right)}{\log_a^2 b+1}  \\
   &  =&\dfrac{3+2 \log_a b-3+2 \log_b a}{\log_a^2 b+1}=\dfrac{2\left(\log_a b+\dfrac{1}{\log_a b}\right)}{\log_a^2 b+1} \\
   & =&\dfrac{2\left(\dfrac{\log_a^2  b+1}{\log_a b}\right)}{\log_a^2 b+1}=\dfrac{2}{\log_a b}=2 \log_ba\neq 6\log_b a.
  \end{eqnarray*} }
  \item $3P=6\log_b a  \neq Q$.
  \item $QP=6\log_ab\cdot 2 \log_b  a=12$.
 \end{itemize}
 }
\end{ex}

% nguồn: Ok_dungsai_lop11_Lan1
\begin{ex}%[1D6V2-3]%[ma:14b2cd4e]
 Biết $ \log 2=a $, $ \log 3=b  $.
 \choiceTF 
   {\True $ \log_4 9=\dfrac{b}{a}$}
  {\True  $\log_6 12=\dfrac{2a+b}{a+b}$}
  {\True $  \log_5 6=\dfrac{a+b}{1-a}$}
  {$ \log_5 {30}=\dfrac{b-1}{1-a}$} 
 \loigiai{ 
  \begin{itemchoice}
   \itemch Đúng. Ta có $ \log_4 9=\dfrac{\log 9}{\log 4}=\dfrac{\log 3^2}{\log  2^2}=\dfrac{2\log 3}{2\log 2}=\dfrac{b}{a}$.
   \itemch Đúng. Ta có $ \log_6 12=\dfrac{\log 12}{\log 6}=\dfrac{\log (2^2\cdot 3)}{\log (2\cdot 3)}=\dfrac{2\log 2+\log  3}{\log 2+\log 3}=\dfrac{2a+b}{a+b}$.
   \itemch Đúng. Ta có $ \log_5 6=\dfrac{\log 6}{\log  5}=\dfrac{\log( 2\cdot 3)}{\log (10: 2)}=\dfrac{\log 2+\log 3}{\log 10 -\log 2}=\dfrac{a+b}{1-a}$.
   \itemch Sai. Ta có $ \log_5 {30}=\log_5 (6\cdot 5)=\log_5 6 + \log_5 5  =\log_5 6 +1$.\\
   Lại có $ \log_5 6=\dfrac{\log  6}{\log 5}=\dfrac{\log( 2\cdot 3)}{\log (10: 2)}=\dfrac{\log 2+\log 3}{\log 10 -\log 2}=\dfrac{a+b}{1-a}$.\\
   Do đó $ \log_5  {30}=\log_5 6 +1=\dfrac{a+b}{1-a}+1=\dfrac{b+1}{1-a}$.
  \end{itemchoice}
  }
\end{ex}

% nguồn: Ok_dungsai_lop11_Lan2
\begin{ex}%[1D6V2-3]%[ma:d9f0c7bc]
	Với mọi số thực  dương $a$, $b$ và $a$, $b\neq1$.
	 \choiceTF
		{\True $\log(ab)=\log  a+\log b$}
		{\True $\log_{a^2}(ab)=\dfrac{1}{2}+\dfrac{1}{2}\log_a  b$}
		{$\log_{3}^{2} a^2=2\log_{3}^{2} a$} 
		{Biết  $\ln\dfrac{a+b}{2}=\dfrac{2\ln a+\ln b}{3}$. Khi đó $a^3+b^3=3\left(8a^2b-ab^2\right)$}
	\loigiai{ 
	\begin{itemchoice}
		\itemch Ta có $a$, $b>0$ và $a$, $b\ne1 \Rightarrow  \log(ab)=\log a+\log b$.
		\itemch Ta có $\log_{a^2}(ab)=\dfrac{1}{2}(\log_a a+\log_a b)=\dfrac{1}{2}(1+\log_a b)=\dfrac{1}{2}+\dfrac{1}{2}\log_a b$. 
		\itemch Ta  có $\log_{3}^{2} a^2=\left(\log_3 a^2\right)^2=\left(2\log_3 a\right)^2=4\log_{3}^{2}a$.
		\itemch  Ta có
		$$\begin{aligned}
		& \ln\dfrac{a+b}{2}=\dfrac{2\ln a+\ln b}{3} \Leftrightarrow 3\ln\dfrac{a+b}{2}=\ln a^2+\ln b  \\ \Leftrightarrow & \ln\left(\dfrac{a+b}{2}\right)^3=\ln\left(a^2b\right) \Leftrightarrow \dfrac{(a+b)^3}{8}=a^2b \\ \Leftrightarrow & a^3+3a^2b+2ab^2+b^3=8a^2b \Leftrightarrow a^3+b^3=5a^2b-3ab^2.
	\end{aligned}$$
\end{itemchoice}
 }
\end{ex}

% nguồn: Ok_dungsai_lop11_Lan2
\begin{ex}%[1D6V2-3]%[ma:55499ce3]
	Cho $a=\log_{2} 5$,  $b=\log_3 5$.
	\choiceTF 
		{\True $b>1$} 
		{$ab=\log_2 3$} 
		{$\log_{12}  5=\dfrac{2}{a+b}$}
		{\True Biết $\log_{24} 250=\dfrac{mab+nb}{pa+qb}$ với $m$, $n$,  $p$, $q\in\mathbb{Z}$. Khi đó $A=mnpq=9$}
	 \loigiai{
	\begin{itemchoice}
		\itemch  Ta có $b=\log_3 5>\log_3 3=1$.
		\itemch Ta có $\log_2 3=\dfrac{\log_5 3}{\log_5 2}=\dfrac{1}{\log_3 4}\cdot\log_2 5=\dfrac{a}{b}$. 
		\itemch Ta có $$\begin{aligned}
		\log_{12} 5 &=\dfrac{1}{\log_{5} 12}=\dfrac{1}{\log_5\left(2^2\cdot3\right)}\\ =& \dfrac{1}{\log_5 2^2+\log_5 3}=\dfrac{1}{\frac{2}{\log_2  5}+\frac{1}{\log_3 5}}=\dfrac{1}{\frac{2}{a}+\frac{1}{b}} \\= &\dfrac{1}{\frac{a+2b}{ab}}=\dfrac{ab}{a+2b}.
	\end{aligned}$$
	\itemch Ta có $$\begin{aligned}
	\log_{24}  250 &=\log_{2^3\cdot3} \left(2\cdot5^3\right)=\log_{2^3\cdot3}2+\log_{2^3\cdot3} \left(5^3\right) \\& =\dfrac{1}{\log_2 \left(2^3\cdot3\right)}+\dfrac{3}{\log_5 \left(2^3\cdot3\right)} \\& =\dfrac{1}{3+\log_2 3}+\dfrac{3}{3\log_5 2+\log_5 3}.
\end{aligned}$$
Mà 
$$\begin{aligned} & \log_2 3=\dfrac{\log_2 5}{\log_3 5}=\dfrac{a}{b}; \\& \log_5 2=\dfrac{1}{\log_2 5}=\dfrac{1}{a}; \\ & \log_5 3=\dfrac{1}{\log_3 5}=\dfrac{1}{b}.
\end{aligned}$$
Khi đó 
$$\log_{24}  250=\dfrac{1}{3+\frac{a}{b}}+\dfrac{3}{\frac{3}{a}+\frac{1}{b}}=\dfrac{1}{\frac{3b+a}{b}}+\dfrac{1}{\frac{3b+a}{ab}}=\dfrac{b}{3b+a}+\dfrac{3ab}{3b+a}=\dfrac{3ab+b}{a+3b}.$$
 Vậy $m=3$, $n=1$, $p=1$, $q=3 \Rightarrow mnpq=9$.
\end{itemchoice}
} 
\end{ex}

% nguồn: Ok_dungsai_lop11_Lan2
\begin{ex}%[1D6V2-3]%[ma:4791bfd7]
	Cho số thực $x$,  $y>0$
	 \choiceTF
		{\True $\log_2 x=\log_4  x^2$}
		{$\log_{\sqrt{2}}+\log_{\frac{1}{3}}\left(2x^2y\right)+\log_4\left(2y^2\right)=1$} 
		 {$\log_3 x+\log_3 x^2+\cdots+\log_3 x^{2023}=2023\cdot2024\log_3 x$}
		{Nếu $x=\log_2 6$,  $y=\log_3 10$. Khi đó $\dfrac{xy+2x-y}{x+2}=\log_2 30$}
	\loigiai{ 
	\begin{itemchoice}
		\itemch Ta có $\log_4 x^2=\log_2^2 (x)2=\log_2 x$  do $x>0$.
		 \itemch Ta có
		$$\begin{aligned}
		\log_{\sqrt{2}}x+\log_{\frac{1}{3}}\left(2x^2y\right)+\log_4\left(2y^2\right) 
		 &=2\log_2 x-\log_2 \left(2x^2y\right)+\dfrac{1}{2}\log_2\left(2y^2\right) \\
		&=\log_2\dfrac{x^2y}{x^2y}-1+\dfrac{1}{2}  \\
		&=-\dfrac{1}{2}. 
	\end{aligned}$$
	\itemch Ta có 
	$$\begin{aligned}
	&\log_3 x+\log_3 x^2+\cdots+\log_3 x^{2\,023} \\ 
	=&\log_3\left(x\cdot x^2\cdots x^{2\,023}\right) \\ 
	 =&\log_3 x^{1+2+\cdots+2\,023}=\dfrac{2\,023\cdot(2\,023+1)}{2}\log_3 x.
\end{aligned}$$
\itemch Ta có $x=\log_2 6$, $y=\log_3 10$. \\ 
 Suy ra
$$\begin{aligned}
 \dfrac{xy+2x-y}{x+2}
&=\dfrac{y(x-1)+2x}{x+2}=\dfrac{\log_3 10\cdot(\log_2 6-1)+2\log_2  6}{\log_2 6+2} \\
&=\dfrac{\log_2 10+\log_2 36}{\log_2 6+2}=\dfrac{\log_2 10+\log_2 36}{\log_2 24}=\log_{24}  360.
\end{aligned}$$
\end{itemchoice}
 }
\end{ex}

% nguồn: Ok_dungsai_lop11_Lan2
\begin{ex}%[1D6V2-3]%[ma:c25d8273]
	Cho hai số thực $a$ và  $b$, với $1<a<b$.
	 \choiceTF
		{$\log_{a^5}b=5\log_a  b$}
		{\True $\log_a  b>1$}
		 {\True $\log_ab^3+\log_{a^2}b^6=\dfrac{6}{\log_b a}$}
		{Nếu $a^2+b^2=7ab$ thì $\log_3(a+b)=1+\log_3  a+\log_3 b$}
	 \loigiai{
	\begin{itemchoice}
		\itemch Ta  có $\log_a b=\dfrac{1}{5}\log_a b$.
		\itemch Ta có  $\log_a b>\dfrac{1}{5}\log_a a=1$.
		\itemch Ta có $\log_ab^3+\log_{a^2}b^6=3\log_ab+3\log_a  b=6\log_a b=\dfrac{6}{\log_b a}$
		\itemch Ta  có $a^2+b^2=7ab \Leftrightarrow a^2+b^2+2ab=9ab \Leftrightarrow (a+b)^2=9ab$. $(\ast)$ \\
		Do $a$, $b>0 \Rightarrow \heva{ & ab>0  \\& a+b>0}$, lấy logarit cơ số $3$ hai vế của $(\ast)$ ta được
		$$\begin{aligned}
		\log_3(a+b)^2=\log_3 9ab 
		& \Leftrightarrow 2\log_3(a+b)=\log_3  9+\log_3 a+\log_3 b \\
		& \Leftrightarrow \log_3 (a+b)=1+\dfrac{1}{2}(\log_3 a+\log_3  b).
	\end{aligned}$$
\end{itemchoice}
} 
\end{ex}

% nguồn: Ok_dungsai_lop11_Lan2
\begin{ex}%[1D6V2-3]%[ma:4590404f]
	Biết $a=\log_{27}  5$, $b=\log_{8} 7$, $c=\log_{2} 3$.
	\choiceTF 
		 {$c>2$}
		 {\True $a\cdot c=\dfrac{1}{3}\log_2 5$}
		{\True $\dfrac{a\cdot  c}{b}=\log_7 5$}
		{\True $\log_{12} 35=\dfrac{3(b+ac)}{c+2}$} 
	 \loigiai{
	\begin{itemchoice}
		\itemch Ta  có $c=\log_2 3<\log_2 4=2$.
		\itemch  Ta có $a\cdot c=c\cdot a=\log_2 3\cdot\log_{27} 5=\log_2 3\cdot\dfrac{1}{3}\log_3 5=\dfrac{1}{3}\log_2 5$.
		\itemch Ta có  $\dfrac{a\cdot c}{b}=\dfrac{\frac{1}{3}\log_2 5}{\log_8 7}=\dfrac{1}{3}\log_2 5\cdot\log_7 8=\dfrac{1}{3}\log_2 5\cdot3\cdot\log_7 2=\log_2 5$.
		\itemch Ta có $\heva{ & a=\log_{27} 5=\dfrac{1}{3}\log_3 5 \\& b=\log_8 7=\dfrac{1}{3}\log_2 7} \Leftrightarrow \heva{& \log_3 5=3a \\&  \log_2 7=3b.}$ \\
		 Mà $\log_{12} 35=\dfrac{\log_2(7\cdot5)}{\log_2\left(3\cdot2^2\right)}=\dfrac{\log_2 7+\log_2 5}{\log_2 3+2}=\dfrac{\log_2 7+\log_2 3\cdot\log_3 5}{\log_2 3+2}=\dfrac{3b+c\cdot3a}{c+2}=\dfrac{3(b+ac)}{c+2}$.
	\end{itemchoice}
	} 
\end{ex}

% nguồn: Ok_dungsai_lop11_Lan2
\begin{ex}%[1D6V2-3]%[ma:b2bef132]
	Cho số thực dương  $a\ne1$.
	 \choiceTF
		{$\ln  a-\ln2=\ln(a-2)$}
		{\True $\log_5\left(125a^2\right)=3+2\log_5  a=0$}
		{\True $2\log_9 a-\dfrac{\log_2  a}{\log_2 3}=0$}
		{$\dfrac{1}{\log_2 a}+\dfrac{1}{\log_3 a}+\cdots+\dfrac{1}\log_{100}a=\log_a 5\,050$} 
	\loigiai{ 
	\begin{itemchoice}
		 \itemch Ta có $\ln a-\ln2=\ln\dfrac{a}{2}$.
		 \itemch Ta có $\log_5\left(125a^2\right)=\log_5 5^3+\log_5a^2=3+2\log_5 a$
		\itemch Ta có $2\log_9 a-\dfrac{\log_2 a}{\log_2 3}=\log_3 a-\log_3 a=0$. 
		 \itemch Ta có
		$$\begin{aligned}
		& \dfrac{1}{\log_2 a}+\dfrac{1}{\log_3 a}+\cdots+\dfrac{1}\log_{100}a \\ =& \log_a 2+\log_a 3+\log_a 4+\cdots+\log_a 99+\log_a 100 \\ =&\log_a  (2\cdot3\cdot4\ldots99\cdot100)=\log_a (100!).
	\end{aligned}$$
\end{itemchoice}
} 
\end{ex}

% nguồn: Ok_dungsai_lop11_Lan2
\begin{ex}%[1D6V2-3]%[ma:c8704d36]
	Cho $\log_{27} 5=a$, $\log_3 7=b$  và $\log_2 3=c$.
	 \choiceTF
		{\True  $0<a<1$}
		 {$bc=\log_7 2$}
		{\True $\log_{6}  45=\dfrac{6ac+2c}{1+c}$}
		{\True  Biết $\log_6 35=\dfrac{(ma+nb)c}{1+pc}$ với $m$, $n$, $p\in\mathbb{Z}$. Khi đó $m^2+n^2+p^2=11$}
	\loigiai{ 
	\begin{itemchoice}
		\itemch Ta có $0=\log_{27} 1<\log_{27} 5<\log_{27}27=1 \Rightarrow  0<a<1$.
		 \itemch Ta có $bc=\log_3 7\cdot\log_2 3=\dfrac{\log_3 7}{\log_3 2}=\log_2 7$.
		\itemch Ta có $\log_{27} 5=a \Rightarrow a=\dfrac{1}{3}\log_3 5\Rightarrow 3a=\log_3  5$. \\
		Khi đó $$\log_{6} 45=\dfrac{\log_2 \left(3^2\cdot5\right)}{\log_2 \left(2\cdot3\right)}=\dfrac{2\log_2 3+\log_2 5}{1+\log_2 3}=\dfrac{2c+\log_2 3\cdot\log_3  5}{1+c}=\dfrac{2c+2c\cdot3a}{1+c}=\dfrac{6ac+2c}{1+c}.$$
		\itemch Ta có $$\log_6 35=\log_6 5+\log_6 7=\dfrac{1}{\log_5 6}+\dfrac{1}{\log_7 6}=\dfrac{1}{\log_5 2+\log_5 3}+\dfrac{1}{\log_7  2+\log_7 3}.$$
		Mà \\ 
		$$\begin{aligned} & \log_{27}5=a \Rightarrow a=\dfrac{1}{3}\log_3 5\Rightarrow 3a\log_3 5 \Rightarrow \log_5 3=\dfrac{1}{3a}; \\& 3ac=\log_3 5\cdot\log_2 3=\log_2 5 \Rightarrow \log_5 2=\dfrac{1}{3ac}; \\ & \log_3 7=b \Rightarrow \log_7 3=\dfrac{1}{b}; \\& bc=\log_2 3\cdot\log_3 7=\log_2 7 \Rightarrow \log_7 2=\dfrac{1}{bc};
	\end{aligned}$$
	Khi  đó
	$$\log_6  35=\dfrac{1}{\dfrac{1}{3ac}+\dfrac{1}{3a}}+\dfrac{1}{\dfrac{1}{b}+\dfrac{1}{bc}}=\dfrac{(3a+b)c}{c+1}.$$
	Vậy $m=3$, $n=1$,  $p=1 \Rightarrow m^2+n^2+p^2=11$.
\end{itemchoice}
} 
\end{ex}

% nguồn: Ok_dungsai_lop11_Lan2
\begin{ex}%[1D6V2-3]%[ma:d46ec036]
	Cho ba điểm $A\left(b;\log_a b\right)$, $B\left(c;2\log_a c\right)$ và $C\left(b;3\log_a b\right)$ với $a$,  $b$, $c$ dương và đều khác $1$.
	\choiceTF 
		{\True $B$ là trung  điểm của $AC$ khi và chỉ khi $b=c > 0$}
		{Ba điểm $A$, $B$, $C$  tạo thành tam giác vuông tại $A$ khi và chỉ khi $b=2c$}
		{Khi ba điểm $A$,  $B$, $C$ tạo thành tam giác thì có diện tích là $S=\left|\left(c-b\right)\cdot\log_a b^2\right|$}
		{\True Khi $B$ là trọng tâm của tam giác $OAC$ (với $O$ là gốc tọa độ) thì giá trị của biểu  thức $S=2b+c$ là bằng $9$}
	 \loigiai{
		\begin{itemchoice}
			 \itemch
			$B$ là trung điểm của $AC$ khi và chỉ khi\\ 
			$\heva{&x_A+x_C=2x_B \\  &y_A+y_C=2y_B
			}\Leftrightarrow \heva{&b+b=2c \\ &\log_a b+3\log_a b=2\left(2\log_a c\right)}\\ \Leftrightarrow \heva{&2b=2c \\ &4\log_a b=4\log_a c} \Leftrightarrow \heva{&b=c \\ &\log_a b=\log_a c}  \Leftrightarrow b=c > 0$.
			\itemch 
			Ta có $\heva{&\overrightarrow{AB}=\left(c-b; 2\log_a c-\log_a b\right)\neq \overrightarrow0 \\&\overrightarrow{AC}=\left(0; 2\log_a b\right)\neq \overrightarrow 0 }$ với $\heva{&\forall b > 0  \\ & b\ne 1.}$\\
			Tam giác $ABC$ vuông ở $A$  $\Leftrightarrow \overrightarrow{AB}\perp \overrightarrow{AB}\Leftrightarrow \overrightarrow{AB}\cdot\overrightarrow{AB}=0$ $\Leftrightarrow 2\log_a b\left(2\log_a c-\log_a b\right)=0$
			$\Leftrightarrow 2\log_a c-\log_a b=0 
			\Leftrightarrow \log_a c^2=\log_a b\Leftrightarrow b=c^2$. 
			\itemch Diện tích tam giác $ABC$ là $S=\dfrac{1}{2} \left|\left(c-b\right)\cdot2\log_a b-0\cdot\left(2\log_a c-\log_a  b\right)\right|$ $=\left|\left(c-b\right)\cdot\log_a b\right|$.
			\itemch 
			Vì $B$  là trọng tâm của tam giác $OAC$ nên\\
			$\heva{&\dfrac{0+b+b}{3}=c \\ &\dfrac{0+\log_a b+3\log_a b}{3}=2\log_a c}\Leftrightarrow \heva{&b+b=3c \\ &4\log_a b=6\log_a c}\Leftrightarrow \heva{&2b=3c \\&2\log_a b=3\log_a c}\\ \Leftrightarrow \heva{&2b=3c \\&\log_a b^2=\log_a c^3}\Leftrightarrow \heva{&2b=3c  \\&b^2=c^3} \Rightarrow b=\dfrac{27}{8}, c=\dfrac{9}{4}.$ \\
			Vậy $S=2b+c=9$. 
		\end{itemchoice}
	 }
\end{ex}


%%%%%%%%%%%%%%%% TRẢ LỜI NGẮN %%%%%%%%%%%%%%%%

% nguồn: Ok_TraLoiNgan_Lop11_Lan1
\begin{ex}%[1D6H2-3]%[ma:8f0291c4]
	Biết $a$, $b$ là các số thực dương thỏa mãn $\log_3 a^2+\log_3  b=5$. Tính giá trị biểu thức $P=5a^2b$.
	 \shortans{1215}
	 \loigiai{
		Ta có $\log_3 a^2+\log_3 b=5\Leftrightarrow \log_3 \left(a^2\cdot b\right)=5\Leftrightarrow  a^2b=243$.\\
		Vậy  $P=5a^2b=5\cdot 243=1215$.
	} 
\end{ex}

% nguồn: Ok_TraLoiNgan_Lop11_Lan1
\begin{ex}%[1D6H2-3]%[ma:f7e94bca]
	Cho $ x,\,y$ là các số thực lớn hơn $  1$ thỏa mãn $x^2+9y^2=6xy$. Tính $$ M=\dfrac{1+\log_{12}x+\log_{12}y}{2\log_{12}\left(x+3y\right)}.$$
	\shortans{1} 
	 \loigiai{
		 Ta có $x^2+9y^2=6xy\Leftrightarrow{\left(x-3y\right)^2}=0\Leftrightarrow x=3y$.\\
	Suy ra $ M=\dfrac{1+\log_{12}x+\log_{12}y}{2\log_{12}\left(x+3y\right)}=  \dfrac{\log_{12}12+\log_{12}3y+\log_{12}y}{2\log_{12}\left(3y+3y\right)}=\dfrac{\log_{12}36y^2}{\log_{12}36y^2}=1$.}
\end{ex}

% nguồn: Ok_TraLoiNgan_Lop11_Lan1
\begin{ex}%[1D6H2-3]%[ma:c32d79a5]
	Cho $\log_a b =3$ và $\log_a c =4$ với $a$,  $b$, $c>0$, $a\ne 1$. Tính giá trị của biểu thức
	$$P=\log_a  \left( \dfrac{a^2\cdot \sqrt{b}}{c^3}\right).$$
	 \shortans{-8,5}
	\loigiai{ 
		\begin{eqnarray*}
			P&=&\log_a \left( \dfrac{a^2\sqrt{b}}{c^3}\right)=\log_a  a^2 +\log_b\sqrt{b}-\log_a c^3\\
			&=&2+\log_a  b^{\frac{1}{2}}-3\log_a c\\
			 &=&2+\dfrac{1}{2}\log_a b -3\log_a c\\ &=&2+\dfrac{3}{2}-12=-\dfrac{17}{2}= -8{,}5.
	\end{eqnarray*}}
\end{ex}

% nguồn: Ok_TraLoiNgan_Lop11_Lan2
\begin{ex}%[1D6H2-3]%[ma:5f7f088a]
	Cho $a$,  $b$, $c$ là các số thực dương $(a, b\ne 1)$ và $\log_ab=5$, $\log_bc=7$.\\
	Tính giá trị của  biểu thức $P=\log_{\sqrt{a}}\left(\dfrac{b}{c}\right)$.
	 \shortans{-60}
	 \loigiai{
		 Ta có
		\begin{align*}
			P&=2\log_a\left(\dfrac{b}{c}\right)\\ 
			 &=2(\log_ab-\log_ac)\\
			 &=2(5-\log_ab\cdot\log_bc)\\
			&=2(5-5\cdot 7)=-60. 
		\end{align*}
	} 
\end{ex}

% nguồn: Ok_TraLoiNgan_Lop11_Lan2
\begin{ex}%[1D6H2-3]%[ma:7ecaf02e]
	Cho $a$, $b$ là  các số thực dương thỏa mãn $\log_ab=3$. Giá trị của biểu thức $I=\log_{\tfrac{a^2}{b}}(ab)$ là
	 \shortans{-4}
	\loigiai{ 
		Ta  có
		\begin{align*}
			I=\log_{\tfrac{a^2}{b}}(ab) =\dfrac{\log_a(ab)}{\log_a\left(\dfrac{a^2}{b}\right)}  =\dfrac{1+\log_ab}{2-\log_ab} =\dfrac{1+3}{2-3}=-4.
		\end{align*}
	} 
\end{ex}

% nguồn: Ok_TraLoiNgan_Lop11_Lan2
\begin{ex}%[1D6H2-3]%[ma:23a0d0da]
	Cho $\log_{700}490=a+\dfrac{b}{c+\log 7}$ với $a,b,c$ là các số nguyên. Tính tổng $T=a+b+c$ 
	\shortans{1} 
	 \loigiai{
		Ta có 
		$$\log_{700}490=\dfrac{\log490}{\log700}  =\dfrac{\log10+\log49}{\log100+\log7} =\dfrac{1+2\log7}{2+\log7} =\dfrac{4+2\log7-3}{2+\log7} =2+\dfrac{-3}{2+\log7}.$$
		Suy ra $a=2$, $b=-3$,  $c=2.$
		Vậy $T=1$. 
	 }
\end{ex}

% nguồn: Ok_TraLoiNgan_Lop11_Lan1
\begin{ex}%[1D6V2-3]%[ma:dce44068]
	Cho $a$, $b$, $x>0$; $a>b$ và $b$, $x\neq 1$ thỏa mãn $\log_x\dfrac{a+2b}{3}=\log_x\sqrt{a}+\dfrac{1}{\log_bx^2}$.  Khi đó biểu thức $P=\dfrac{2a^2+3ab+b^2}{(a+2b)^2}$ có giá trị bằng
	\shortans{1,25} 
	\loigiai{	Ta có \allowdisplaybreaks\begin{eqnarray*}
		&&\log_x\dfrac{a+2b}{3}=\log_x\sqrt{a}+\dfrac{1}{\log_bx^2}\\ 
		&\Leftrightarrow&\log_x\dfrac{a+2b}{3}=\log_x\sqrt{a}+\log_x\sqrt{b} \\ 
		 &\Leftrightarrow& a+2b=3\sqrt{ab}\\
		&\Leftrightarrow&  a^2-5ab+4b^2=0\\
		&\Leftrightarrow&(a-b)(a-4b)=0\Leftrightarrow a=4b. 
	\end{eqnarray*}
Vậy $P=\dfrac{2a^2+3ab+b^2}{(a+2b)^2}=\dfrac{32b^2+12b^2+b^2}{36b^2}=\dfrac{5}{4}=1{,}25$.} 
\end{ex}

% nguồn: Ok_TraLoiNgan_Lop11_Lan2
\begin{ex}%[1D6V2-3]%[ma:665f37f5]
	Cho $\log_23=a$  và $\log_35=b$. Biết $\log_{12}150=\dfrac{2ab+ma+n}{a+2}$ với $m,\, n\in \mathbb{Z}$. Giá trị của $L=m+n$ là
	\shortans{2} 
	\loigiai{ 
		Ta có $\log_23=a  \Leftrightarrow \log_32 =\dfrac{1}{a}$ và $\log_35=b$.\\
		Khi  đó
		\begin{align*}
			\log_{12}150=\dfrac{\log_3150}{\log_312} =\dfrac{\log_3(2\cdot3\cdot5^2)}{\log_3(2^2\cdot3)} =\dfrac{\log_32+\log_33+2\log_35}{2\log_32+\log_33} 
			=\dfrac{\dfrac{1}{a}+1+2b}{\dfrac{2}{a}+1} 
			 =\dfrac{2ab+a+1}{a+2}.
		\end{align*}
		 Do $m,\, n\in \mathbb{Z}$ $\Rightarrow m=1$ và $n=1$ $\Rightarrow L=m+n=2$.
	 }
\end{ex}

% nguồn: Ok_TraLoiNgan_Lop11_Lan2
\begin{ex}%[1D6V2-3]%[ma:a14ff302]
	Cho $x=2024!$. Giá trị của biểu thức $A=\dfrac{1}{\log_2x}+\dfrac{1}{\log_3x}+\cdots+\dfrac{1}{\log_{2024}x}$  là
	\shortans{1} 
	 \loigiai{
		Theo bài $x=2024! \Rightarrow  x>0$, $x\ne 1$.\\
		\begin{eqnarray*}\allowdisplaybreaks
			 A&=&\dfrac{1}{\log_2x}+\dfrac{1}{\log_3x}+\cdots+\dfrac{1}{\log_{2024}x} \\
			 &=&\log_x2+\log_x3+\cdots+\log_x2024 \\
			&=&\log_x(1\cdot 2\cdot 3\cdot  2024) \\
			 &=&\log_x2024!.
		\end{eqnarray*}
		Với $x=2024!\Rightarrow A=\log_{2024!}2024!=1$. 
	 }
\end{ex}

% nguồn: Ok_TraLoiNgan_Lop11_Lan2
\begin{ex}%[1D6V2-3]%[ma:e05d309f]
	Cho $0<a\ne 1$. Tính giá trị của biểu  thức $B=\log_a\left(\dfrac{a^2\cdot\sqrt[3]{a}\cdot\sqrt[5]{a^4}}{\sqrt[4]{a}}\right) +a^{2\log_a\tfrac{\sqrt{105}}{30}}$
	\shortans{3} 
	 \loigiai{
		Ta  có
		\begin{align*}
			 B&=\log_a\left(\dfrac{a^2\cdot\sqrt[3]{a}\cdot\sqrt[5]{a^4}}{\sqrt[4]{a}}\right) +a^{2\log_a\tfrac{\sqrt{105}}{30}}\\
			&=\log_a\left(\dfrac{a^2\cdot a^{\tfrac{1}{3}}\cdot a^{\tfrac{4}{5}}}{a^{\tfrac{1}{4}}}\right) +\left(a^{\log_a\tfrac{\sqrt{105}}{30}}\right)^2\\ 
			 &=\log_aa^{\tfrac{173}{60}}+\left(\dfrac{\sqrt{105}}{30}\right)^2\\
			&=\dfrac{173}{60}+\dfrac{105}{900}=3. 
		\end{align*}
	 }
\end{ex}

% nguồn: Ok_TraLoiNgan_Lop11_Lan2
\begin{ex}%[1D6V2-3]%[ma:2cdd9c80]
	Cho $a$, $b$, $c$ là ba số thực  dương, khác $1$, thỏa mãn $\log_a b=2\log_b c=4\log_c a$ và $a+2b+3c=48$. Tính $P=a+b+c$.
	\shortans{21} 
	 \loigiai{
		Do $a$, $b$, $c$ đều khác  $1$ nên $\log_a b$, $\log_b c$ và $\log_c a$ đều khác $0$. \\
		Ta có 
		$$ 
		\log _a  b=2 \log _b c
		\Leftrightarrow \log _a c \cdot \log _c b=2 \log _b c 
		\Leftrightarrow \log _a c=2 \log _b^2 c 
		 $$
		và 
		 $$
		\log _a b=4 \log _c a \Leftrightarrow \log _a c \cdot \log _c b=4 \log _c a \Leftrightarrow \log _c b=4 \log _c^2 a. 
		$$ 
		Suy  ra
		$$ 
		\log  _a c \cdot \log _c b=8 \log _b^2 c \cdot \log _c^2 a
		\Leftrightarrow \log _a  b=8 \log _b^2 a
		\Leftrightarrow \log _a  b=2
		\Leftrightarrow  b=a^2.
		 $$
		Ta có $\log _a b=2 \log _b  c=2 \Rightarrow b=c$.\\
		Theo giả  thiết
		$$ 
		a+2b+3c=48 
		\Leftrightarrow a+2a^2+3a^2=48 
		\Leftrightarrow  5a^2+a-48=0
		\Leftrightarrow \hoac{&a=-\dfrac{16}{5} \\  &a=3..}
		$$ 
		Do các số $a$, $b$, $c$ dương  nên $a=3 \Rightarrow b=c=9$. \\
		 Vậy, $P=a+b+c=21$.
	} 
\end{ex}

% nguồn: Ok_TraLoiNgan_Lop11_Lan2
\begin{ex}%[1D6V2-3]%[ma:bf1ad690]
	Cho $\log_{6} 45=a+\dfrac{\log_{2} 5+b}{\log_{2} 3+c}$, với  $a, b, c \in \mathbb{Z}$. Tính tổng $a+b+c$.
	 \shortans{1}
	 \loigiai{
		Ta  có
		\begin{align*}
			 \log_{6} 45
			&=\log_{6}(3^2  \cdot 5)
			=\dfrac{\log_{2}(3^2 \cdot 5)  }{\log_{2} 6}
			 =\dfrac{2\log_{2} 3+\log_{2} 5}{\log_{2}(2 \cdot 3) }\\
			&=\dfrac{2\log_{2} 3+\log_{2} 5}{1+\log_{2} 3} 
			=\dfrac{2\left(\log_{2}  3+1\right)+\log_{2} 5-2}{\log_{2} 3+1}\\
			&=2+\dfrac{\log_{2} 5-2}{\log_{2}  3+1}.
		\end{align*}
		Vậy $\heva{&a=2 \\ &b=-2  \\ &c=1} \Rightarrow a+b+c=2-2+1=1$.
	} 
\end{ex}

% nguồn: Ok_TraLoiNgan_Lop11_Lan2
\begin{ex}%[1D6V2-3]%[ma:1012151e]
	Gọi $n$ là số nguyên dương sao cho:$\dfrac{1}{\log_{2023} x}+\dfrac{1}{\log_{2023^2} x}+\dfrac{1}{\log_{2023^3} x}+ \hdots +\dfrac{1}{\log_{2023^n} x}=\dfrac{210}{\log_{2023} x}$ đúng với  mọi $x$ dương, $x \neq 1$. Tính giá trị của biểu thức $P=3n+4$.
	 \shortans{64}
	\loigiai{ 
		Với mọi $x$ dương, $x  \neq 1$. \\
		Ta có 
		\begin{align*}
			\dfrac{1}{\log_{2023} x} + \dfrac{1}{\log_{2023^2} x} + \hdots + \dfrac{1}{\log_{2023^n}  x}
			&= \log_{x} 2023  + \log_{x} 2023^2 + \log_{x} 2023^3 + \hdots + \log_{x} 2023^n \\
			 &=(1+2+3+ \hdots +n) \log_{x} 2023 \\
			 &= \dfrac{n(n+1) }{2}\log_{x} 2023
		\end{align*}
		 Do đó theo yêu cầu đề bài, ta có
		$$ 
		\dfrac{n(n+1) }{2} =  210
		\Leftrightarrow  n^2+n-420=0
		\Leftrightarrow  \hoac{&n=20 \\ &n=-21.}
		 $$
	Do đó $n  = 20$. Khi đó $P = 3n+4 = 64$.}
\end{ex}

% nguồn: Ok_TraLoiNgan_Lop11_Lan2
\begin{ex}%[1D6V2-3]%[ma:3709923d]
	Cho biểu thức $P=\dfrac{1}{\log_{a^2} 2}+\dfrac{1}{\log_{a^4} 2}+\ldots+\dfrac{1}{\log_{a^{2n}} 2}$, với $a$ là số thực lớn hơn $1$, $n \in \mathbb{N}^*$. Biết $P \le  100 \log_2 a$, tính tổng các giá trị của $n$.
	 \shortans{45}
	\loigiai{ 
		 Ta có
		\begin{align*}
			P &=  \dfrac{1}{\log_{a^2} 2} + \dfrac{1}{\log_{a^4} 2} + \ldots + \dfrac{1}{\log_{a^{2n}}2}\\
			&= \log_2 a^2 + \log_2  a^4 + \ldots + \log_2 a^{2n}\\
			&= \log_2 a^{2+4+\ldots+2n} \\ 
			&= (2+4+\ldots+2n) \log_2 a\\ 
			&= \dfrac{2(1+n) n}{2}\log_2 a \\ 
			&= n(n+1) \log_2 a 
		\end{align*}
		Lại  có
		\begin{align*}
			P \le 100 \log_2 a 
			&\Leftrightarrow n(n+1)  \log_2 a \le 100\log_2 a \\
			&\Leftrightarrow n^2+n-100 \le 0\\ 
			&\Leftrightarrow \dfrac{-1-\sqrt{401}}{2} \le n \le \dfrac{-1+\sqrt{401}}{2} 
		\end{align*}
		Do $n \in \mathbb{N}^*$ nên $n \in \{1;2;3;4;5;6;7;8;9\}$. \\ 
	Vậy, tổng các giá trị của $n$ là  $45$.}
\end{ex}


%%%%%%%%%%%%%%%% TỰ LUẬN %%%%%%%%%%%%%%%%

% nguồn: Ok_TuanLuan_Toan11_Lan1
\begin{bt}%[1D6H2-3]%[ma:0d8c0624]
	Cho $a=\log_25$, $b=\log_35$. Hãy biểu  diễn $\log_65$ theo $a$ và $b$.
	 \loigiai{
		Ta  có $a=\log_25\Rightarrow \log_52=\dfrac{1}{a}$, $b=\log_35\Rightarrow \log_53=\dfrac{1}{b}$.\\
		 Từ đó
		 \[\log_65=\dfrac{1}{\log_56}=\dfrac{1}{\log_52+\log_53}=\dfrac{1}{\dfrac{1}{a}+\dfrac{1}{b}}=\dfrac{ab}{a+b}.\]
	 }
\end{bt}

% nguồn: Ok_TuanLuan_Toan11_Lan1
\begin{bt}%[1D6H2-3]%[ma:459959be]
	Cho các số dương $a,b$  và $a\neq 1$. \\Rút gọn biểu thức $K=\log_ab^2+2\log_{a^2}{b^4}+3\log_{a^3}{b^6}-4\log_{a^4}{b^8}$.
	\loigiai{ 
		\begin{eqnarray*}
			K  &=& \log_ab^2+2\log_{a^2}{b^4}+3\log_{a^3}{b^6}-4\log_{a^4}{b^8}\\
			&=&  2\log_ab+4\log_ab+6\log_ab-8\log_ab\\
			&=&  4\log_ab.
		\end{eqnarray*}
	 }
\end{bt}

% nguồn: Ok_TuanLuan_Toan11_Lan2
\begin{ex}%[1D6H2-3]%[ma:5a1a1dda]
	Cho $a=\log_25$, $b=\log_35$. Hãy biểu diễn  $\log_65$ theo $a$ và $b$.
	\loigiai{ 
		 Ta có $a=\log_25\Rightarrow \log_52=\dfrac{1}{a}$, $b=\log_35\Rightarrow \log_53=\dfrac{1}{b}$.\\
		 Từ đó
		\[\log_65=\dfrac{1}{\log_56}=\dfrac{1}{\log_52+\log_53}=\dfrac{1}{\dfrac{1}{a}+\dfrac{1}{b}}=\dfrac{ab}{a+b}.\] 
	} 
\end{ex}

% nguồn: Ok_TuanLuan_Toan11_Lan2
\begin{ex}%[1D6H2-3]%[ma:85b409fe]
	Cho  các số dương $a,b$ và $a\neq 1$. \\Rút gọn biểu thức $K=\log_ab^2+2\log_{a^2}{b^4}+3\log_{a^3}{b^6}-4\log_{a^4}{b^8}$.
	 \loigiai{
		\begin{eqnarray*}
			K  &=& \log_ab^2+2\log_{a^2}{b^4}+3\log_{a^3}{b^6}-4\log_{a^4}{b^8}\\
			 &=& 2\log_ab+4\log_ab+6\log_ab-8\log_ab\\
			&=&  4\log_ab.
		\end{eqnarray*}
	 }
\end{ex}
